\appendix

\section*{Methodology Details}
\label{appendix:methodology}

\noindent\textbf{Platform Selection and Categorization.}
Platform selection employed purposive sampling targeting: (1) platforms with documented deployment in commercial VASP environments or significant market presence based on industry reports and regulatory filings; (2) comprehensive coverage across functional categories (transaction monitoring, risk scoring, forensics, Travel Rule, privacy-preserving); (3) diversity in architectural approaches (SaaS, on-premise, hybrid, decentralized); and (4) representation of both established market leaders and recent entrants (post-2020). Academic research was identified through systematic search across IEEE Xplore, ACM Digital Library, arXiv, and Google Scholar using query combinations spanning blockchain technology, compliance mechanisms, and analytical techniques, with inclusion criteria requiring peer-reviewed publication or verifiable prototype implementation.

\noindent\textbf{Compliance Mechanism Selection Criteria.}
The compliance mechanisms summarized in \autoref{tab:compliance_comparison} prioritize concepts that are (1) widely used or explicitly named in industry practice, and (2) operationally distinct rather than purely re-labeled variants of existing mechanisms. Terms that are not recognized as industry-standard (e.g., speculative "Know-Your" labels, like Know Your Model and Know Your Intent), or that substantially overlap with existing mechanisms such as KYT and address/wallet screening, were excluded to avoid taxonomy inflation. This selection focuses the systematization on mechanisms with demonstrable adoption and clear operational scope.

\noindent\textbf{Capability Assessment Methodology.}
Capability assessments in Table~\ref{tab:capability_matrix} are based on analysis of platform documentation, technical whitepapers, published API specifications, and vendor capability statements available as of December 2024. Full support (\fullcirc) indicates production-ready features explicitly documented and commercially available; partial support (\halfcirc) indicates beta/limited features, announced roadmap items, or capabilities restricted to specific blockchain networks; no support (\emptycirc) indicates absence of documented capability. Academic prototypes were assessed based on published research papers and available implementations. These assessments represent capabilities as claimed or documented; independent validation of all features was not feasible. Coverage percentages reflect the proportion of platform categories demonstrating at least partial support for each capability within our analyzed sample.

\section*{Research Limitations}
\label{appendix:limitations}

This analysis is subject to several methodological constraints that readers should consider when interpreting our findings:

\noindent\textbf{Vendor Claim Verification.} Platform capabilities are assessed based on vendor documentation and public claims without independent technical validation of all features, introducing potential overstatement bias. While we cross-referenced claims with academic literature and industry reports where possible, exhaustive feature verification was beyond the scope of this systematization.

\noindent\textbf{Selection Bias.} Selection criteria favor platforms with English-language documentation and public accessibility, potentially underrepresenting solutions in non-English markets or serving exclusively domestic jurisdictions. Platforms requiring enterprise contracts for documentation access were excluded, which may bias our sample toward more transparent vendors.

\noindent\textbf{Market Coverage.} Market share and adoption figures represent estimates based on available industry reports and vendor disclosures rather than comprehensive market census. The rapidly evolving nature of Web3 RegTech means capabilities and platform status may change between data collection (October-December 2025) and publication.

\noindent\textbf{Generalizability.} Quantitative claims (e.g., percentage distributions, coverage statistics) should be interpreted as characterizing our analyzed sample rather than representing comprehensive market statistics, as exhaustive enumeration of all global RegTech solutions is infeasible. Our sample of 41 platforms represents a substantial portion of publicly documented solutions but cannot claim complete market coverage.

\noindent\textbf{Categorical Boundaries.} Categorical classifications represent analytical constructs that may not align with platforms' self-descriptions, and many platforms span multiple categories. The taxonomies presented reflect our systematization framework rather than industry-standard definitions, which largely do not exist for this emerging domain.

\noindent\textbf{Temporal Constraints.} The Web3 RegTech ecosystem evolves rapidly, with new platforms launching and existing platforms adding features continuously. Our analysis represents a snapshot as of late 2025 and may not reflect subsequent developments.

\noindent\textbf{Academic Research Coverage.} While our systematic literature review covered major academic databases and venues, the decentralized nature of blockchain research (with significant contributions appearing in arXiv preprints, workshop papers, and technical reports) means some relevant work may have been missed.
